\section{Autonomous response-memory closures}
\label{sec:closure}

\subsection{Shifted Legendre coordinates}

Let $p_k$ be the shifted Legendre polynomial of degree $k$ on $[0,1]$,
normalized by
\begin{equation}
 p_k(1)=1,qquad
 \int_0^1p_j(x)p_k(x)\,\od x=\frac{\mathbf 1_{j=k}}{2k+1}.
 \label{eq:legendre-normalization}
\end{equation}
Their dilation identity is
\begin{equation}
 x p_k'(x)=k p_k(x)+\sum_{j<k}(2j+1)p_j(x).
 \label{eq:legendre-dilation}
\end{equation}
For order $P$, let $p=(p_0,\ldots,p_{P-1})^\top$, let
$e=(1,\ldots,1)^\top\in\R^P$, and define the lower-triangular matrix
\begin{equation}
 A_{kk}=k,qquad A_{kj}=2j+1\ (j<k),qquad A_{kj}=0\ (j>k),
 \label{eq:A-matrix}
\end{equation}
with indices starting at zero.  Then $xp'(x)=Ap(x)$.

The same identity that makes online Legendre projection possible in signal
memory \citep{gu2020hippo} makes our neural closure autonomous.  When the
history interval expands, differentiating a moment only mixes it with moments
of the same or lower degree.  No stored past samples are needed.

\subsection{Variant A: residual-activity clock}
\label{sec:old-clock}

Set
\begin{equation}
 \dot\tau=\rho,qquad \tau(0)=1.
 \label{eq:activity-clock}
\end{equation}
The unit prefix $0\le\xi\le1$ carries constant forward history
$h_{\ell-1,a}(0)$ and zero backward history.  Physical training inserts the
current pair $(h_{\ell-1,a},b_{\ell,a})$ at
$\xi=\tau(t)$.  Since $\od\xi=\rho\od t$, the learned integral in
\eqref{eq:intro-history} is unchanged.

For every internal link $2\le\ell\le L$ and sample $a$, store $n\times P$
moment matrices $H_{\ell,a}$ and $U_{\ell,a}$ whose $k$th columns are
\begin{align}
H_{\ell,a,:,k}&=\int_0^\tau
 \overline h_{\ell-1,a}(\xi)p_k(\xi/\tau)\,\od\xi,\nonumber\\
U_{\ell,a,:,k}&=\int_0^\tau
 \overline b_{\ell,a}(\xi)p_k(\xi/\tau)\,\od\xi .
\label{eq:old-moments}
\end{align}
Differentiation and \eqref{eq:legendre-dilation} give the closed ODE
\begin{align}
\dot H_{\ell,a}&=\rho h_{\ell-1,a}e^\top
             -\frac{\rho}{\tau}H_{\ell,a}A^\top,\nonumber\\
\dot U_{\ell,a}&=r_a\delta_{\ell,a}e^\top
             -\frac{\rho}{\tau}U_{\ell,a}A^\top.
\label{eq:old-moment-ode}
\end{align}
Initially, the zeroth column of $H_{\ell,a}$ is $h_{\ell-1,a}(0)$ and
all other columns of $H_{\ell,a}$ and $U_{\ell,a}$ vanish.

Orthogonal projection on $[0,\tau]$ has diagonal Gram matrix
$G_A=\tau\,\mathrm{diag}((2k+1)^{-1})$.  The reconstructed physical matrix is
\begin{equation}
 \boxed{
 \what W_\ell=W_{\ell,0}
 -\frac{2}{nm\tau}\sum_{a=1}^m\sum_{k=0}^{P-1}(2k+1)
 U_{\ell,a,:,k}H_{\ell,a,:,k}^\top .}
 \label{eq:old-reconstruction}
\end{equation}
All forward and backward responses in \eqref{eq:old-moment-ode} are evaluated
using the current reconstructed matrices.  The first layer and readout follow
their canonical dense equations evaluated at the same state.  Equations
\eqref{eq:model-forward}--\eqref{eq:old-reconstruction} therefore form an
autonomous finite ODE.  They do not use the future or the dense trajectory.

For one sample and one internal link, the lowest moments make the meaning
transparent:
\begin{equation}
 H_{:,0}(t)=h(0)+\int_0^t\rho(s)h(s)\,\od s,qquad
 U_{:,0}(t)=\int_0^t r(s)\delta(s)\,\od s.
\end{equation}
Higher modes distinguish where and how the responses changed along accumulated
activity.  They are history projections, not Taylor derivatives at
initialization.

The learned correction in \eqref{eq:old-reconstruction} has rank at most
$mP$ per internal layer.  The action on a vector $v$ is evaluated without
forming a dense learned matrix:
\begin{equation}
 (\what W_\ell-W_{\ell,0})v
 =-\frac{2}{m\tau}\sum_{a,k}\frac{2k+1}{n}
 U_{\ell,a,:,k}\bigl(H_{\ell,a,:,k}^\top v\bigr).
 \label{eq:factor-action}
\end{equation}
The action of $W_{\ell,0}$ and the corresponding transpose action remain
present exactly.

\subsection{Variant B: response-arclength clock}
\label{sec:new-clock}

The activity clock controls the residual factor but does not directly control
the speed of every stored response.  Concatenate all histories that enter the
learned internal matrices:
\begin{equation}
 \Psi=\operatorname{concat}_{\ell=2}^L
       \operatorname{concat}_{a=1}^m
       \bigl(h_{\ell-1,a},b_{\ell,a}\bigr).
 \label{eq:Psi}
\end{equation}
On an interval with $\rho>0$, define
\begin{equation}
 \boxed{\dot\tau=g:=\rho+\norm{\dot\Psi}_2,qquad \tau(0)=1.}
 \label{eq:response-clock}
\end{equation}
This is a causal function of the current closure state.  Its key property is
\begin{equation}
 \left\lVert\frac{\od\Psi}{\od\tau}\right\rVert_2
 =\frac{\norm{\dot\Psi}_2}{\rho+\norm{\dot\Psi}_2}\le1.
 \label{eq:unit-speed}
\end{equation}
The clock controls the joint path speed; it does not equalize every neuron and
is not claimed to minimize finite-order projection error.

Changing coordinates must not change the learning integral.  We therefore
advance the location of history at speed $g$ but insert mass at speed $\rho$.
For a test function $v$ define
\begin{equation}
 \int v(\xi)\,\od\mu_t(\xi)
 =\int_0^1v(\xi)\,\od\xi
  +\int_0^t v(\tau(s))\rho(s)\,\od s .
 \label{eq:weighted-measure}
\end{equation}
On the physical-history portion, $\od\mu_t/\od\xi=\rho/g\le1$.
Both prefix histories now match their initial values:
$h(\xi)=h(0)$ and $b(\xi)=b(0)$ for $0\le\xi\le1$.  The known prefix product
is subtracted from the reconstruction so it causes no artificial learning.

Define weighted moments and one shared $P\times P$ Gram matrix,
\begin{align}
H_{\ell,a}&=\int h_{\ell-1,a}(\xi)p(\xi/\tau)^\top\,\od\mu_t(\xi),\nonumber\\
U_{\ell,a}&=\int b_{\ell,a}(\xi)p(\xi/\tau)^\top\,\od\mu_t(\xi),\nonumber\\
G&=\int p(\xi/\tau)p(\xi/\tau)^\top\,\od\mu_t(\xi).
\label{eq:weighted-moments}
\end{align}
They satisfy
\begin{align}
\dot H_{\ell,a}&=\rho h_{\ell-1,a}e^\top
                   -\frac g\tau H_{\ell,a}A^\top,\nonumber\\
\dot U_{\ell,a}&=r_a\delta_{\ell,a}e^\top
                   -\frac g\tau U_{\ell,a}A^\top,\nonumber\\
\dot G&=\rho ee^\top-\frac g\tau(AG+GA^\top).
\label{eq:weighted-ode}
\end{align}
At initialization, only the zeroth columns of $H$ and $U$ are nonzero and
$G=\mathrm{diag}((2k+1)^{-1})$.  The reconstruction is
\begin{equation}
 \boxed{
 \what W_\ell=W_{\ell,0}-\frac{2}{nm}\sum_{a=1}^m
 \left[U_{\ell,a}G^{-1}H_{\ell,a}^\top
       -b_{\ell,a}(0)h_{\ell-1,a}(0)^\top\right].}
 \label{eq:weighted-reconstruction}
\end{equation}
The prefix makes $G$ positive definite at every fixed finite order and finite
clock length.  In computation one solves linear systems with $G$ rather than
forming $G^{-1}$.  The correction rank is at most $m(P+1)$.

\subsection{Why the new clock is explicit}
\label{sec:explicit-clock}

At first sight \eqref{eq:response-clock} may appear implicit because
$\dot\Psi$ depends on matrix velocities.  Differentiating
\eqref{eq:weighted-reconstruction} resolves the issue.  Set
\begin{equation}
 h_{\ell-1,a}^\star=H_{\ell,a}G^{-1}e,qquad
 b_{\ell,a}^\star=U_{\ell,a}G^{-1}e.
\end{equation}
Every dilation term proportional to $g/\tau$ cancels, leaving
\begin{equation}
 V_\ell:=\dot{\what W}_\ell
 =-\frac{2\rho}{nm}\sum_a\left[
 b_{\ell,a}(h_{\ell-1,a}^\star)^\top
 +b_{\ell,a}^\star(h_{\ell-1,a}-h_{\ell-1,a}^\star)^\top
 \right].
 \label{eq:physical-velocity}
\end{equation}
This velocity is known before $g$.  Set $V_1=F_1(\what\theta)$ and
$V_w=F_w(\what\theta)$ and differentiate the forward pass:
\begin{align}
\dot z_{1,a}&=V_1u_a,\nonumber\\
\dot z_{\ell,a}&=V_\ell h_{\ell-1,a}
                    +\what W_\ell\dot h_{\ell-1,a},\nonumber\\
\dot h_{\ell,a}&=\phi_\ell'(z_{\ell,a})\odot\dot z_{\ell,a},\nonumber\\
\dot r_a&=\frac1n\bigl(V_w^\top h_{L,a}+w^\top\dot h_{L,a}\bigr),
\qquad
\dot\rho=\frac1{m\rho}\sum_a r_a\dot r_a.
\label{eq:directional-forward}
\end{align}
Differentiate backpropagation from layer $L$ down to layer two:
\begin{align}
\dot\delta_{L,a}
 &=V_w\odot\phi_L'(z_{L,a})
  +w\odot\phi_L''(z_{L,a})\odot\dot z_{L,a},\nonumber\\
\dot\delta_{\ell,a}
 &=\phi_\ell''(z_{\ell,a})\odot\dot z_{\ell,a}
       \odot\what W_{\ell+1}^\top\delta_{\ell+1,a}\nonumber\\
 &\quad+\phi_\ell'(z_{\ell,a})\odot
  \left(V_{\ell+1}^\top\delta_{\ell+1,a}
       +\what W_{\ell+1}^\top\dot\delta_{\ell+1,a}\right),\nonumber\\
\dot b_{\ell,a}
 &=\frac{\dot r_a\delta_{\ell,a}+r_a\dot\delta_{\ell,a}}{\rho}
   -b_{\ell,a}\frac{\dot\rho}{\rho}.
\label{eq:directional-backward}
\end{align}
These formulas determine $\dot\Psi$, then $g$, and finally
\eqref{eq:weighted-ode}.  The algorithm uses directional derivatives, fixed
matrix actions, moment contractions, and a small Gram solve.  It uses no full
Jacobian, Hessian, future trajectory, or dense learned matrix.  Notice that
$\phi_1''$ is unnecessary: the differentiated backward histories begin at
layer two.

If $\rho(0)=0$, both dense and closure trajectories are stationary; we define
the closure to remain there without evaluating $r/\rho$.  The theorem below
also shows that a nonstationary regular trajectory cannot reach zero residual
at finite time under its hypotheses.

